\section{Introduction}
In this chapter, we consider the problem of modeling a time series of several variables gathered at irregularly-spaced times, especially when only a few variables are typically measured at one time. Our motivating application is to model the sequence of vital signs and laboratory values recorded over time in the EHR of a hospitalized patient. Effective models of such data could enable many improvements to both patient care and hospital logistics, such as extrapolating the future state of the patient, facilitating diagnosis of underlying illness, or crafting personalized treatment strategies~\citep{hyland2020early, bica2021real}. 

Realizing these benefits requires a model that can handle the two key aspects of our chosen scenario. First, the model must accommodate the \emph{irregular timing of individual measurements}. Second, the model must capture the \emph{data-dependent patterns of missingness}. To clarify this latter point, whether a specific variable is measured or not in the present moment is often influenced by current or recent measurements of other variables. Intuitively, doctors rarely order lab tests unrelated to the patient's present condition or medical history.

We focus on a specific technical formulation of the generative mechanism for missingness, known as \emph{missing at random} (MAR) \citep{rubin1976inference}. In MAR settings, the presence or absence of a measurement variable (binary indicator of missingness) depends on the values of all observed (non-missing) variables. As one example of MAR, the current value of Troponin, an indicator of heart failure measured via lab processing of a blood draw, could be missing due to other observed vitals indicating normal heart function. As another example, the absence of cholesterol measurements could be attributed to the lack of evident risk factors for heart disease, alongside a normal weight and healthy Body Mass Index (BMI). In contrast to MAR, the  commonly assumed \emph{missing completely at random} (MCAR) setting assumes the missingness indicator arises independently of any observations.

Unfortunately, many common models for time-varying labs and vitals are not capable of handling both irregular timing and data-driven (MAR) missingness assumptions.
Well-known recurrent neural nets for clinical time-series, such as GRU-D~\citep{che2018recurrent} or BRITS~\citep{caoBRITSBidirectionalRecurrent2018}, do allow for arbitrary missingness at each timestep at both training and prediction.
However, these methods assume regular time intervals, necessitating alignment to a uniformly-spaced time grid as a pre-processing step.
%Furthermore, these methods assume MCAR rather than MAR mechanisms.
Another growing line of work directly models irregularly timed measurements via continuous formulation of recurrent networks \citep{pineda1987generalization, pearlmutter1989learning} and more recently via neural networks coupled with differential equation solvers \citep{de2019gru, rubanova2019latent}. However, in these works (as well as works like GRU-D and BRITS) there is no focus on a data-dependent (MAR) missingness mechanism. The current landscape thus lacks a method that can simultaneously handle irregular-spacing and data-dependent missingness.

In this work, our \textbf{primary contribution} is a method that fills this gap. We present a probabilistic generative model for an irregular time series of clinical measurements with an data-driven mechanism that allows arbitrary missingness at any time. 
Irregularly-spaced measurements are handled via an encoder-decoder stochastic differential equation architecture known as the continuous recurrent unit (CRU)~\citep{schirmer2022modeling}. Data driven missingness is handled via regression on time-specific representation vectors produced by the CRU. Our approach  preserves the efficient closed-form updates inside the CRU.

We demonstrate the effectiveness of our approach by examining prediction error
when \emph{extrapolating} the next few hours of a patient's multivariate time series. Across 3 real health records datasets, each containing thousands of patients, we show our approach delivers lower extrapolation error (averaged across 12+ labs and vitals) than several state-of-the-art alternative methods designed for irregularity but not data-dependent missingness. Further sensitivity analyses on synthetic data confirm that our advantage occurs specifically in settings where the MCAR assumptions of baselines are not appropriate.
While our evaluations focus on EHRs for hospitalized patients, our method design is general purpose and useful for any time-series modeling that needs both irregular time and data-dependent missingness.

As a \textbf{contribution to best practices in evaluation}, our experiments report two kinds of error. First, we report
average error across all channels after normalization (as many past works do). Second, we report channel-specific error without normalization, using the raw units of the original physiological measurement. This latter approach is not routinely done in other works, but we find it is valuable way to contextualize our method's gains.


\subsection*{Generalizable Insights about Machine Learning in the Context of Healthcare}

Our work highlights the importance of careful joint modeling of missingness and measured values, using a data-driven (MAR) mechanism, while also accounting for the irregularly-spaced nature of EHR records over time. The choice of what to measure and when carries critical information that can be used to improve understanding of the patient's condition. Hopefully this improved understanding and ability to forecast future patient state could help lead to improved treatment decisions and better outcomes.
